Since the pendulum relation uses SI units, each length is first converted from millimetres to metres. For set \(j\), the small-angle relation then gives

\[\boxed{\omega_{n,j}=\sqrt{\frac{g}{L_j}}}\qquad\text{(A.2.1.a)}\]

The frequency and period follow from \(\omega_{n,j}=2\pi f_{n,j}\) and \(f_{n,j}=1/\tau_j\), so

\[\boxed{f_{n,j}=\frac{\omega_{n,j}}{2\pi}}\qquad\text{(A.2.1.b)}\]

\[\boxed{\tau_j=\frac{2\pi}{\omega_{n,j}}}\qquad\text{(A.2.1.c)}\]

Using \(g=9.81\ \mathrm{m/s^2}\) and the ten lengths gives the following values for A.2.1.d.

{ % do not increment counter
\begin{center}\small\begin{tabular}{@{}rrrrr@{}}
\toprule\noalign{}
Set & \(L\) (mm) & \(\omega_n\) (rad/s) & \(f_n\) (Hz) & \(\tau\) (s) \\
\midrule\noalign{}
1 & 502.013 & 4.42 & 0.704 & 1.421 \\
2 & 499.337 & 4.43 & 0.705 & 1.418 \\
3 & 500.472 & 4.43 & 0.705 & 1.419 \\
4 & 502.949 & 4.42 & 0.703 & 1.423 \\
5 & 502.216 & 4.42 & 0.703 & 1.422 \\
6 & 496.635 & 4.44 & 0.707 & 1.414 \\
7 & 500.416 & 4.43 & 0.705 & 1.419 \\
8 & 498.255 & 4.44 & 0.706 & 1.416 \\
9 & 498.350 & 4.44 & 0.706 & 1.416 \\
10 & 499.358 & 4.43 & 0.705 & 1.418 \\
\bottomrule\end{tabular}\end{center}
}

The mean of each calculated column follows from \(\overline{q}=\frac{1}{10}\sum_{j=1}^{10}q_j\). Therefore, the values required in A.2.2 are

\[\boxed{\overline{\omega_n}=4.43\ \mathrm{rad/s}}\qquad\text{(A.2.2)}\]

\[\boxed{\overline{f_n}=0.705\ \mathrm{Hz}}\qquad\text{(A.2.2)}\]

\[\boxed{\overline{\tau}=1.419\ \mathrm{s}}\qquad\text{(A.2.2)}\]

For A.2.3, the initial velocity is zero, so the undamped free response reduces to

\[\boxed{u(t)=u_0\cos(\overline{\omega_n}t)=4\cos(4.429470397t)\ \mathrm{mm}}\qquad\text{(A.2.3)}\]

At \(t_1=18\ \mathrm{s}\), substitution into the response gives

\[\boxed{u(t_1)=4\cos\!\left(4.429470397\times18\right)\ \mathrm{mm}=-1.484\ \mathrm{mm}}\qquad\text{(A.2.3.a)}\]

For the requested fifteen-cycle duration, multiply the mean period by \(N_{\mathrm{osc}}=15\). The duration and corresponding position are

\[\boxed{t_{N_{\mathrm{osc}}}=15\overline{\tau}=21.28\ \mathrm{s}}\qquad\text{(A.2.3.b)}\]

\[\boxed{u_{N_{\mathrm{osc}}}=u(15\overline{\tau})=4.00\ \mathrm{mm}}\qquad\text{(A.2.3.b)}\]

The mean of the individual periods is \(\overline{\tau}=1.4185009419\ \mathrm{s}\), whereas the period associated with the mean frequency is \(2\pi/\overline{\omega_n}=1.4184958345\ \mathrm{s}\). Averaging and taking a reciprocal do not commute. Following the prompt, the duration uses \(15\overline{\tau}\); the difference from fifteen exact cycles of the mean-frequency model is only \(76.6\ \mu\mathrm{s}\), and the endpoint still rounds to \(4.00\ \mathrm{mm}\).

Using the boxed A.2.3 expression over \(0\leq t\leq15\overline{\tau}\) gives the fifteen-oscillation response below. Since the marker is placed at \(t_1=18\ \mathrm{s}\), the plotted value is

\[\boxed{u(18)=-1.484\ \mathrm{mm}}\qquad\text{(A.2.3.c)}\]

\begin{figure}[H]
\centering
\includegraphics[width=0.95\linewidth,height=\textheight,keepaspectratio,alt={A.2.3.c. Pendulum response over fifteen oscillations with a data tip at eighteen seconds.}]{figures/generated/solutions/A2.pdf}
\caption{A.2.3.c. Pendulum response over fifteen oscillations with a data tip at \ensuremath{t=18\,\mathrm{s}}.}
\end{figure}

\clearpage
\matlabcaption{Script used to reproduce the A.2 response figure.}

\lstinputlisting[style=hwmatlab]{matlab/A2.m}
